% ------------------------- Capítulo II -----------------------------
% Definições
\newtheorem{def:base-va}{Definição}[chapter]
\newtheorem{def:va}[def:base-va]{Definição}
\newtheorem{def:espaço induzido}[def:base-va]{Definição}
\newtheorem{def:evento va}[def:base-va]{Definição}
\newtheorem{def:cdf}[def:base-va]{Definição}
\newtheorem{def:suporte}[def:base-va]{Definição}
\newtheorem{def:variavel discreta}[def:base-va]{Definição}
\newtheorem{def:fmp}[def:base-va]{Definição}
\newtheorem{def:variavel continua}[def:base-va]{Definição}
\newtheorem{def:fdp}[def:base-va]{Definição}
\newtheorem{def:transformacao}[def:base-va]{Definição}
\newtheorem{def:variavel derivada}[def:base-va]{Definição}
\newtheorem{def:esperanca}[def:base-va]{Definição}
\newtheorem{def:momento}[def:base-va]{Definição}
\newtheorem{def:momentos importantes}[def:base-va]{Definição}
\newtheorem{def:mgf}[def:base-va]{Definição}
\newtheorem{def:simetria}[def:base-va]{Definição}
\newtheorem{def:mediana}[def:base-va]{Definição}
\newtheorem{def:moda}[def:base-va]{Definição}
% \newtheorem{def:independencia de variaveis}[def:base-va]{Definição}
% \newtheorem{def:covariancia}[def:base-va]{Definição}

% Proposições
\newtheorem{prop:base-va}{Proposição}[chapter]
\newtheorem{prop:suporte contem imagem}[prop:base-va]{Proposição}
\newtheorem{prop:intersecao suporte}[prop:base-va]{Proposição}
\newtheorem{prop:prob fora suporte}[prop:base-va]{Proposição}
\newtheorem{prop:ep discreto}[prop:base-va]{Proposição}
\newtheorem{prop:steltjes e fda}[prop:base-va]{Proposição}
\newtheorem{prop:suporte enumeravel}[prop:base-va]{Proposição}
\newtheorem{prop:fmp}[prop:base-va]{Proposição}
\newtheorem{prop:fda discreto}[prop:base-va]{Proposição}
\newtheorem{prop:existencia fmp}[prop:base-va]{Proposição}
\newtheorem{prop:existencia fdp}[prop:base-va]{Proposição}
\newtheorem{prop:fdp e fda}[prop:base-va]{Proposição}
\newtheorem{prop:continuidade fda}[prop:base-va]{Proposição}
\newtheorem{prop:densidade 1}[prop:base-va]{Proposição}
\newtheorem{prop:densidade fora suporte}[prop:base-va]{Proposição}
\newtheorem{prop:prob transformacao}[prop:base-va]{Proposição}
\newtheorem{prop:padronizacao}[prop:base-va]{Proposição}
\newtheorem{prop:mgf da transformacao afim}[prop:base-va]{Proposição}
\newtheorem{prop:equivalencia simetria}[prop:base-va]{Proposição}
\newtheorem{prop:simetria e mediana}[prop:base-va]{Proposição}
% \newtheorem{prop:variancia da soma}[prop:base-va]{Proposição}

% Lema
\newtheorem{lem:base-va}{Lema}[chapter]
\newtheorem{lem:prob conjunto enumeravel}[lem:base-va]{Lema}
\newtheorem{lem:densidade continua}[lem:base-va]{Lema}

% Teoremas
\newtheorem{thm:base-va}{Teorema}[chapter]
\newtheorem{thm:cdf}[thm:base-va]{Teorema}
\newtheorem{thm:prob continua}[thm:base-va]{Teorema}
\newtheorem{thm:prob suporte}[thm:base-va]{Teorema}
\newtheorem{thm:suporte continua}[thm:base-va]{Teorema}
\newtheorem{thm:prob transformacao discreta}[thm:base-va]{Teorema}
\newtheorem{thm:fda transformacao continua}[thm:base-va]{Teorema}
\newtheorem{thm:fdp transformacao continua}[thm:base-va]{Teorema}
\newtheorem{thm:fdp transformacao nao monotona continua}[thm:base-va]{Teorema}
\newtheorem{thm:esperanca de transformacao}[thm:base-va]{Teorema}
\newtheorem{thm:propriedades esperanca}[thm:base-va]{Teorema}
\newtheorem{thm:propriedades variancia}[thm:base-va]{Teorema}
\newtheorem{thm:derivada do momento}[thm:base-va]{Teorema}
\newtheorem{thm:condicao de igualdade de distribuicao e momento}[thm:base-va]{Teorema}
\newtheorem{thm:convergencia mgf}[thm:base-va]{Teorema}
\newtheorem{thm:simetria funcao de probabilidade}[thm:base-va]{Teorema}
\newtheorem{thm:simetria media moda mediana}[thm:base-va]{Teorema}
% \newtheorem{thm:propriedades covariancia}[thm:base-va]{Teorema}
% \newtheorem{thm:covariancia independentes}[thm:base-va]{Teorema}

% Corolário
\newtheorem{cor:base-va}{Corolário}[chapter]
\newtheorem{cor:cdf}[cor:base-va]{Corolário}
\newtheorem{cor:simetria e media aritmetica}[cor:base-va]{Corolário}

% Observações
\newtheorem{obs:base-va}{Observação}[chapter]
\newtheorem{obs:simetria 1}[obs:base-va]{Observação}
\newtheorem{obs:simetria 2}[obs:base-va]{Observação}
